\section{Risk alignment through mask geometry}
\label{sec:theory}

\noindent\textbf{Flagship instance: full surface Hausdorff.}
Let \(d_{H,x}\) be the full Hausdorff distance between inner digital surfaces,
normalized by the image diagonal and extended with zero empty--empty cost and
unit one-sided-empty cost. It is a metric on the resulting surface-set space
and a pseudometric on filled masks. The deployment cost is
\(L_x^{\mathrm{HD}}(A,B)=d_{H,x}(A,B)\in[0,1]\). The reverse triangle
inequality gives
\begin{equation*}
  |L_x^{\mathrm{HD}}(Y,A)-L_x^{\mathrm{HD}}(Y',A)|
  \leq d_{H,x}(Y,Y'),
\end{equation*}
and the general transport lemma in \Cref{app:theory} therefore yields
\begin{equation*}
  \Delta_{L^{\mathrm{HD}},\Yhat}(x)
  \leq W_{1,d_{H,x}}(P_x,Q_{p(x)})
  \leq \TV(P_x,Q_{p(x)}).
\end{equation*}
This is the \textbf{main geometry-aware guarantee}. Under \(Q_p\), one scalar level moves
the candidate surface coherently through the probability landscape; under HD,
the cost records the most displaced point of that surface. The construction and
deployment risk are thus aligned at the level of boundary motion. Combining this display
with \Cref{sec:framework} controls HD score error and then HD AURC
regret.

\noindent\textbf{Robust extension: pooled HD95.}
Replacing the maximum pooled surface distance by its 95th percentile suppresses
a small extreme tail. The resulting normalized penalized HD95 cost remains in
\([0,1]\), so the threshold integral, total-error decomposition, scalar cost-law
transport bound, and AURC theorem all apply. It is not a mask metric, however:
percentile costs can jump when the fraction of distant surface points crosses
five percent and need not satisfy a triangle inequality. Consequently HD95 is
a \textbf{practical robust extension} of the boundary experiment, not a corollary of the
constant-one HD Wasserstein theorem. \Cref{app:theory} gives an explicit
quantile-jump construction.

\noindent\textbf{Regional contrast: Dice.}
For binary masks, define \(\Dice(A,B)=2|A\cap B|/(|A|+|B|)\),
with \(\Dice(\varnothing,\varnothing)=1\) and Dice zero when exactly one mask
is empty. The deployment cost \(L_{\Dice}=1-\Dice\) lies in \([0,1]\). For Jaccard
distance \(d_J\), \Cref{app:theory} proves
\begin{equation*}
  \Delta_{L_{\Dice},\Yhat}(x)
  \leq\min\{\TV(P_x,Q_{p(x)}),2W_{1,d_J}(P_x,Q_{p(x)})\}.
\end{equation*}
Thus Dice and HD instantiate the \textbf{same Lipschitz--Wasserstein template} under
different, noninterchangeable mask geometries. Dice does not, however, make the
comonotone level-set coupling uniquely natural in the way coordinated surface
motion motivates it for HD.

\Cref{eq:quad} estimates expected Dice under \(Q_p\). Because level
sets change only at probability knots, sorting tied groups gives the exact
integral in \(O(N\log N)\) time. The algorithm in \Cref{app:theory}
removes Dice quadrature error; midpoint integration remains the common estimator
used for controlled cross-risk comparison.

\noindent\textbf{Why SDC and level-set Dice can be close.}
When its denominator is positive, SDC is
\begin{equation}
  \SDC(p,\Yhat)
  =\frac{2\sum_i p_i\hat y_i}{\sum_i p_i+\sum_i\hat y_i},
  \qquad \hat y_i=\mathbf 1\{i\in\Yhat\}.
  \label{eq:sdc}
\end{equation}
For any mask law \(Q\) with marginals \(\Ex_QY_i=p_i\), let
\(N(Y)=2\sum_iY_i\hat y_i\) and
\(D(Y)=\sum_iY_i+\sum_i\hat y_i\). Then
\(\SDC=\Ex_QN/\Ex_QD\), whereas posterior Dice is \(\Ex_Q[N/D]\); they are
not generally equal. For a nonempty action of size \(m\), however,
\begin{equation*}
  \left|\Ex_Q L_{\Dice}(Y,\Yhat)-[1-\SDC(p,\Yhat)]\right|
  \leq \frac{2}{m}\sqrt{\Var_Q |Y\cap\Yhat|}
       +\frac{1}{2m}\sqrt{\Var_Q |Y\setminus\Yhat|}.
\end{equation*}
SDC is therefore the \textbf{count-mean plug-in approximation}: spatial coupling affects
posterior Dice only through the induced overlap/outside-count law, and its
deviation from SDC is controlled by count dispersion. Published SDC theory
connects the ratio to a product-Bernoulli target under exact marginals and
conditional pixel independence \citep{sdc}; our level-set score is exact under
the different, comonotone \(Q_p\). Neither assumption dominates the other.

\noindent\textbf{A marginal-preserving coupling extension.}
To test whether Dice needs non-nested regional variation, we embed the global
level-set law in a spatial Gaussian-copula family. For independent standard
normal \(G,\varepsilon_i\) and a unit-variance spatial Gaussian field \(U\),
we sample
\begin{equation*}
  Z_i=\sqrt{\alpha}G+\sqrt{\beta}U_i+\sqrt{1-\alpha-\beta}\,\varepsilon_i,
  \qquad Y_i=\mathbf 1\{Z_i\leq\Phi^{-1}(p_i)\}.
\end{equation*}
Every \(Z_i\) is standard normal, hence \(\Pr(Y_i=1)=p_i\) exactly. The family
connects the global shared-threshold endpoint \((\alpha,\beta)=(1,0)\), smooth
non-nested masks, and independent Bernoulli sampling at \((0,0)\). We implement
\(U\) as a standardized bilinear coordinate field, avoiding a dense covariance
matrix. This is a \textbf{coupling-sensitivity experiment}, not a recovered true
posterior: it introduces spatial-scale choices and Monte Carlo error, while the
global method remains parameter-free and exactly integrable.
